\section*{Bab \ref{ch:g12:meiosis-genetic-shuffling} --- Meiosis dan Pengocokan Genetik}

\begin{solution}{exo:g12:meiosis-genetic-shuffling:1}
\omterm{def:g12:meiosis-genetic-shuffling:meiosis}{Diploid}: dua salinan tiap \omterm{def:g11:cell-cycle-mitosis:chromosome}{kromosomnya} ($2n$); \omterm{def:g12:meiosis-genetic-shuffling:meiosis}{haploid}: satu ($n$).
\omterm{def:g12:meiosis-genetic-shuffling:meiosis}{Meiosis} membagi dua jumlah \omterm{def:g11:cell-cycle-mitosis:chromosome}{kromosomnya}, dari $2n$ menjadi $n$, dan
menurunkan \omterm{def:g10:universal-dna:information}{DNA-nya} dari $2q$ (setelah replikasi) menjadi $q/2$ tiap gamet.
\end{solution}

\begin{solution}{exo:g12:meiosis-genetic-shuffling:2}
\omterm{def:g12:meiosis-genetic-shuffling:meiosis}{Meiosis} I memisahkan kromosom homolog tiap pasangannya; \omterm{def:g12:meiosis-genetic-shuffling:meiosis}{meiosis} II
memisahkan kedua \omterm{def:g11:cell-cycle-mitosis:chromosome}{kromatid} tiap \omterm{def:g11:cell-cycle-mitosis:chromosome}{kromosomnya}. Yang pertama bersifat reduksi.
\end{solution}

\begin{solution}{exo:g12:meiosis-genetic-shuffling:3}
Setelah \omterm{def:g12:meiosis-genetic-shuffling:meiosis}{meiosis} I: 4 \omterm{def:g11:cell-cycle-mitosis:chromosome}{kromosom} (masing-masing dua \omterm{def:g11:cell-cycle-mitosis:chromosome}{kromatid}), \omterm{def:g10:universal-dna:information}{DNA} $q$.
Setelah \omterm{def:g12:meiosis-genetic-shuffling:meiosis}{meiosis} II: 4 \omterm{def:g11:cell-cycle-mitosis:chromosome}{kromosom} (masing-masing satu \omterm{def:g11:cell-cycle-mitosis:chromosome}{kromatid}), \omterm{def:g10:universal-dna:information}{DNA} $q/2$.
\end{solution}

\begin{solution}{exo:g12:meiosis-genetic-shuffling:4}
$2^4 = 16$.
\end{solution}

\begin{solution}{exo:g12:meiosis-genetic-shuffling:5}
Persilangan dengan pasangan yang hanya membawa \omterm{def:g10:universal-dna:gene}{alel} \omterm{def:g11:genes-and-disease:genetic}{resesif}, sehingga
tiap keturunannya memperlihatkan gamet tetua yang diuji. Bagian yang
rekombinan menyingkapkan apakah dua \omterm{def:g10:universal-dna:gene}{gennya} berada pada \omterm{def:g11:cell-cycle-mitosis:chromosome}{kromosom} yang
sama dan, jika ya, seberapa jauh jaraknya.
\end{solution}

\begin{solution}{exo:g12:meiosis-genetic-shuffling:6}
\omterm{def:g12:meiosis-genetic-shuffling:meiosis}{Meiosis} I: $2q$; di antara pembelahannya: $q$; gametnya: $q/2$. Tanpa
fase S di antaranya, pembelahan keduanya membagi dua \omterm{def:g10:universal-dna:information}{DNA-nya} lagi, dan
itulah yang membawa gametnya menjadi separuh isi \omterm{def:g10:cells-common-unit:cell}{sel} \omterm{def:g12:meiosis-genetic-shuffling:meiosis}{diploid}.
\end{solution}

\begin{solution}{exo:g12:meiosis-genetic-shuffling:7}
Empat kelas yang sama besar: tidak terpaut, berada pada \omterm{def:g11:cell-cycle-mitosis:chromosome}{kromosom} yang
berbeda. Gametnya $AB$, $Ab$, $aB$, $ab$, masing-masing seperempat.
\end{solution}

\begin{solution}{exo:g12:meiosis-genetic-shuffling:8}
Terpaut: dua kelas besar ($AB$, $ab$, tetua) dan dua kelas kecil ($Ab$,
$aB$, rekombinan). Rekombinannya $85/1000 = 8.5\%$.
\end{solution}

\begin{solution}{exo:g12:meiosis-genetic-shuffling:9}
Hanya $AB$ dan $ab$: pertukaran di atas kedua \omterm{def:g10:universal-dna:gene}{gennya} menukar ruas yang
tidak membawa satu pun \omterm{def:g10:universal-dna:gene}{gennya}, sehingga gabungan $A/a$ dan $B/b$-nya
tidak berubah. Penggabungan ulang antara dua \omterm{def:g10:universal-dna:gene}{gen} menuntut sebuah
pertukaran di antara keduanya.
\end{solution}

\begin{solution}{exo:g12:meiosis-genetic-shuffling:10}
Sekitar 0.8, 2.9 dan 16 tiap 1000: faktornya 20 antara 25 dan 42.
\end{solution}

\begin{solution}{exo:g12:meiosis-genetic-shuffling:11}
Pada \omterm{def:g12:meiosis-genetic-shuffling:meiosis}{meiosis} I, kedua homolognya pergi ke satu \omterm{def:g10:cells-common-unit:cell}{sel}: dua gamet dengan
dua salinan dan dua gamet tanpa salinan. Pada \omterm{def:g12:meiosis-genetic-shuffling:meiosis}{meiosis} II, kedua
\omterm{def:g11:cell-cycle-mitosis:chromosome}{kromatid} satu \omterm{def:g11:cell-cycle-mitosis:chromosome}{kromosomnya} pergi ke satu \omterm{def:g10:cells-common-unit:cell}{sel}: satu gamet dengan dua
salinan, satu tanpa salinan, dan dua yang normal.
\end{solution}

\begin{solution}{exo:g12:meiosis-genetic-shuffling:12}
Penggabungan ulangnya sebanding dengan jaraknya: 2\% berarti \omterm{def:g10:universal-dna:gene}{gennya}
sangat berdekatan, sehingga pertukarannya jarang jatuh di antaranya;
48\% berarti berjauhan, dengan pertukaran di antaranya pada nyaris tiap
\omterm{def:g12:meiosis-genetic-shuffling:meiosis}{meiosisnya}. Karena satu pertukaran memberi 50\% gamet rekombinan, \omterm{def:g10:universal-dna:gene}{gen}
terpaut yang berjauhan bergabung ulang sebebas \omterm{def:g10:universal-dna:gene}{gen} yang bebas.
\end{solution}

\begin{solution}{exo:g12:meiosis-genetic-shuffling:13}
$AB$, $Ab$, $aB$, $ab$, masing-masing seperempat, lewat pemilahan bebas
kedua pasangannya. \omterm{prop:g12:meiosis-genetic-shuffling:crossingover}{Pindah silang} tidak mengubah apa pun bagi \omterm{def:g10:universal-dna:gene}{gen} ini:
ia menggabungkan ulang \omterm{def:g10:universal-dna:gene}{alel} di sepanjang satu \omterm{def:g11:cell-cycle-mitosis:chromosome}{kromosom}, dan \omterm{def:g10:universal-dna:gene}{gen} ini
berada pada \omterm{def:g11:cell-cycle-mitosis:chromosome}{kromosom} yang berbeda.
\end{solution}

\begin{solution}{exo:g12:meiosis-genetic-shuffling:14}
Tiap orang tuanya mewariskan kepada tiap anaknya salah satu dari dua
\omterm{def:g10:universal-dna:gene}{alelnya} pada tiap \omterm{def:g10:universal-dna:gene}{gennya}, yang terpilih secara acak; pada sebuah \omterm{def:g10:universal-dna:gene}{gen}
tertentu, dua bersaudara menerima \omterm{def:g10:universal-dna:gene}{alel} yang sama dari seorang tua
dengan peluang $1/2$. Sepanjang banyak \omterm{def:g10:universal-dna:gene}{gen}, bagian yang dibaginya
rata-rata $1/2$, tetapi bagi sepasang saudara tertentu ia berubah-ubah
di sekitarnya.
\end{solution}

\begin{solution}{exo:g12:meiosis-genetic-shuffling:15}
Klon tumbuhan induknya, yang secara genetik serupa dengannya dan serupa
satu sama lain. Tumbuhannya memperoleh pelipatgandaan yang cepat dan
pasti bagi \omterm{def:g11:enzymes-and-phenotype:phenotype}{genotipe} yang bekerja di sini dan sekarang; ia kehilangan
keragaman yang akan membiarkan sebagian keturunannya selamat dari
perubahan lingkungan atau parasit yang baru.
\end{solution}

\begin{solution}{pb:g12:meiosis-genetic-shuffling:1}
\textbf{1.} Betinanya $b^+b$, $vg^+vg$, dengan $b^+vg^+$ pada satu
\omterm{def:g11:cell-cycle-mitosis:chromosome}{kromosom} dan $b\,vg$ pada \omterm{def:g11:cell-cycle-mitosis:chromosome}{kromosom} yang lain (dari kedua orang tuanya
yang murni). Jantannya $bb$, $vg\,vg$.

\textbf{2.} Hanya gamet $b\,vg$: jantannya hanya menyumbangkan \omterm{def:g10:universal-dna:gene}{alel}
\omterm{def:g11:genes-and-disease:genetic}{resesif}, sehingga \omterm{def:g11:enzymes-and-phenotype:phenotype}{fenotipe} tiap keturunannya menyingkapkan gamet betinanya.

\textbf{3.} $b^+vg^+$ (kelabu-normal), $b\,vg$ (hitam-vestigial),
$b^+vg$ (kelabu-vestigial), $b\,vg^+$ (hitam-normal).

\textbf{4.} Kelabu-normal 42\%, hitam-vestigial 41\%, kelabu-vestigial
9\%, hitam-normal 8\%.

\textbf{5.} Pada \omterm{def:g11:cell-cycle-mitosis:chromosome}{kromosom} yang sama: dua kelas besar dan dua kelas
kecil alih-alih $1:1:1:1$.

\textbf{6.} Kelabu-normal dan hitam-vestigial: yaitu gabungan yang
dibawa kedua homolognya, yang diwarisi utuh dari orang tuanya yang murni.

\textbf{7.} Kelabu-vestigial dan hitam-normal, dari sebuah \omterm{prop:g12:meiosis-genetic-shuffling:crossingover}{pindah
silang} antara kedua \omterm{def:g10:universal-dna:gene}{gennya} pada profase I.

\textbf{8.} $(206 + 185)/2300 \approx 17\%$. Sebuah \omterm{prop:g12:meiosis-genetic-shuffling:crossingover}{pindah silang}
antara \omterm{def:g10:universal-dna:gene}{gennya} membuat dua dari keempat \omterm{def:g11:cell-cycle-mitosis:chromosome}{kromatidnya} menjadi rekombinan,
sehingga 17\% gamet rekombinan berarti sebuah pertukaran pada sekitar
34\% \omterm{def:g12:meiosis-genetic-shuffling:meiosis}{meiosisnya}.

\textbf{9.} Satu pertukaran menghasilkan kedua \omterm{def:g11:cell-cycle-mitosis:chromosome}{kromatid} rekombinannya
bersamaan, masing-masing satu macam; dan kedua homolognya pergi ke
gametnya sama rata, sehingga kedua kelas tetuanya juga sepadan.

\textbf{10.} Kelas tetuanya kini kelabu-vestigial dan hitam-normal,
masing-masing sekitar 41.5\%; rekombinannya kelabu-normal dan
hitam-vestigial, masing-masing sekitar 8.5\%.

\textbf{11.} Ya: penggabungan ulangnya jauh di bawah 50\% dengan keduanya.

\textbf{12.} $b$ --- $cn$ --- $vg$: 9 dan 8 berjumlah 17.

\textbf{13.} Peluang sebuah pertukaran dalam sebuah selang sebanding
dengan panjangnya, sehingga frekuensi selang yang bersebelahan
bertambah: frekuensi penggabungan ulangnya mengukur jarak di sepanjang
\omterm{def:g11:cell-cycle-mitosis:chromosome}{kromosomnya}.

\textbf{14.} Satu pertukaran memberi dua \omterm{def:g11:cell-cycle-mitosis:chromosome}{kromatid} rekombinan dari
empat: paling banyak 50\%; pertukaran tambahannya mengocok ulang tetapi
tidak pernah menaikkan bagian rekombinannya di atas separuh.

\textbf{15.} \omterm{def:g10:universal-dna:gene}{Gen} itu berada pada \omterm{def:g11:cell-cycle-mitosis:chromosome}{kromosom} lain atau sangat jauh pada
\omterm{def:g11:cell-cycle-mitosis:chromosome}{kromosom} yang sama; persilangannya tidak dapat memberitahukan yang mana.

\textbf{16.} $2^4 = 16$ gamet; $16 \times 16 = 256$ gabungan.

\textbf{17.} Tiap pertukarannya menghasilkan \omterm{def:g11:cell-cycle-mitosis:chromosome}{kromosom} yang tidak pernah
ada sebelumnya, pada kedudukan yang berubah dari satu \omterm{def:g12:meiosis-genetic-shuffling:meiosis}{meiosis} ke
\omterm{def:g12:meiosis-genetic-shuffling:meiosis}{meiosis} berikutnya; jumlah \omterm{def:g11:cell-cycle-mitosis:chromosome}{kromosom} yang berbeda, dan dengan demikian
jumlah gametnya, tidak punya batas dalam praktiknya.

\textbf{18.} $2^{23} \times 2^{23} \approx \num{7e13}$: tujuh ratus
kali lebih banyak gabungan daripada jumlah manusia yang pernah hidup,
dan itu sebelum \omterm{prop:g12:meiosis-genetic-shuffling:crossingover}{pindah silangnya} dihitung.

\textbf{19.} Bersaudara mengundi \omterm{def:g10:universal-dna:gene}{alelnya} dari empat perangkat yang sama
(dua tiap orang tuanya) dan berbagi, rata-rata, separuhnya; orang asing
mengundi dari perangkat yang berbeda.

\textbf{20.} Angka 17\%, yaitu bagian gamet betinanya yang tergabung
ulang antara $b$ dan $vg$, mengukur jarak antara \omterm{def:g10:universal-dna:gene}{gennya}; pemilahan
bebas pasangannya dan \omterm{prop:g12:meiosis-genetic-shuffling:crossingover}{pindah silang} di dalam tiap pasangannya membuat
tiap gametnya berbeda.
\end{solution}
